\section{Conclusion}
\label{sec:conclusion}

Spark ships an automatic remedy for skewed joins, and practitioners continue to
salt by hand. We set out to determine which they should use, and how much salt
to use when salting. The answer on our workload was blunter than expected:
Spark's skew-join rule never engaged at all, so there was nothing to choose
between.

The measurement is a null result and we report it as one. Enabling AQE on a
join whose longest task ran \BaseMaxTaskThick~s produced a longest task of
\AqeCoalMaxTaskThick~s, with speedups of \AqeSpeedupMid$\times$ and
\AqeSpeedupThick$\times$ against measurement floors of 0.1\% and 6.5\%. A
paired control with coalescing disabled shows that what little AQE contributed
was partition merging, which lowers fixed overhead and does nothing for a
straggler. Selective salting, on the same workload, cut the critical path
\SaltTaskCutThick$\times$.

The cause is a precondition that is documented and almost never discussed. A
skewed partition is split into pieces of \texttt{advisoryPartitionSizeInBytes};
a hot partition smaller than that size cannot be split at all, and the rule
becomes a no-op however extreme the skew. Our hot partition measured
\MeasHotMB~MB against a 16~MB advisory size. Sweeping the two parameters whose
names contain the word \emph{skew} to their most permissive settings changed
nothing; lowering the advisory size engaged the rule immediately. Practitioners
tune the two knobs that are named for the problem, and the one that gates the
solution at moderate data volumes is not among them.

On the cost model we are more circumspect than we expected to be. The
replication term is confirmed exactly in shuffle records --- \SelRecSlope{}
rows per salt value for selective salting against \UnifRecSlope{} for uniform,
a ratio of \RecSlopeRatio{} matching the predicted $M/P$ with $R^{2}=1.000$.
Its consequence for runtime is not confirmed: the replication coefficient is
statistically unidentified in all \NumFits{} fits, and in \NumNegativeB{} of
them it comes out with the wrong sign. At this scale replication cost is not
merely unresolved in wall-clock; it is absent. We therefore present the closed
form $k^{*}=\sqrt{\gamma H/P}$ as theory, state plainly that this study does
not validate it, and release the cluster configuration under which the
prediction --- that network shuffle raises the coefficient and lowers $k^{*}$
--- becomes testable.

The decision between mechanisms turns on a statistic that costs one sampled
group-by to obtain. When skew is one-sided and above AQE's trigger, AQE alone
suffices and manual salting adds replication for little gain. When the same key
is heavy on both sides, split-and-replicate degenerates and salting remains
necessary. When the hot partition sits below AQE's absolute threshold, the rule
does nothing at all --- silently --- and a practitioner who has enabled AQE
believes the problem is handled.

The artifact matters as much as the result. We cannot out-scale the
organisations whose systems papers describe this machinery at production
volume, and we do not try. What a practitioner-built harness can contribute is
specification: a public, reproducible generator with a declared skew taxonomy,
a metrics layer that separates deterministic counters from machine-dependent
timings, and an honest account of what a single node does and does not license.
\textsc{skewbench} is released under the MIT licence.

Three directions follow. The most useful is validation at cluster scale, where
the network shuffle makes $b$ larger and $k^{*}$ correspondingly smaller than
we report; our prediction is directional and testable. The second is to fold
$\gamma$ estimation into the engine, so that $k^{*}$ is computed rather than
configured --- salting is currently the one skew remedy Spark cannot apply for
you, and the model here is simple enough to be a planner rule. The third is to
extend the taxonomy: multi-key joins, outer joins, and skew that drifts over
time are all common in production and none is covered here.
